\section{Learning Skills and Composing the Task}\label{sec:method}
\subsection{Adapting the Action Interface}
The target setting is a new embodiment for which the pretrained action interface cannot be used directly. The required intervention can range from a new output projection or controller head to replacement of an action expert. We retain useful VLA representations and learn the action-producing components needed by the new controller. This distinction matters: creating a residual action module does not imply that the complete pretrained expert was reinitialized.

In the evaluated SmolVLA implementation, the pretrained policy, including its action expert and projections, is loaded and retained. Its predicted actions are mapped to the target interface by coordinate selection and normalization. A separate residual MLP is initialized from scratch for initial skill learning; subsequent phases inherit its learned weights. Action-expert LoRA adds trainable low-rank updates in earlier phases while preserving the pretrained base weights. The final continuation freezes the inherited VLA/LoRA and trains the residual and exploration parameters. Thus, this experiment evaluates a newly learned action module over a retained VLA, not a reset of the complete pretrained expert.

\subsection{Dense Rewards for Initial Skill Learning}
We divide one object-manipulation task into skills such as approaching the object, picking it up, and transporting and releasing it at the destination. Each skill is initially trained with a dense reward reflecting progress toward its own goal. For example, approach feedback can depend on reduction of gripper-to-object distance, while pickup feedback can distinguish grasp formation and lifting. This gives early rollouts informative feedback even before binary success occurs.

Skill training uses its declared reset distribution and instruction. The task ``move to the apple'' therefore supplies approach experience; it is not counted as a successful execution of ``put the apple on the plate.'' Similarly, a held-object placement episode teaches transport and release without claiming that the policy performed the preceding grasp. \pending{Insert the dense reward functions, skill resets, and training budgets for the selected lineage.}

\subsection{Binary Refinement and Self-Imitation}
After initial skill acquisition, training transitions to a binary reward indicating whether the instructed skill was completed. This shifts the learning criterion from partial progress to successful execution. We collect successful policy-generated trajectories and use them as supervised targets, retaining experience from earlier skills when consolidating a new one. Here self-imitation means fitting successful experience generated by the policy or a retained policy, not an additional human-demonstration requirement.

For successful skill banks $\mathcal D_k^+$, a shared consolidation objective has the form
\begin{equation}
\mathcal L_{\rm SI}(\theta)=\sum_{k=1}^{K}\alpha_k
\mathbb E_{(o,\ell,a)\sim\mathcal D_k^+}
\big[\ell_{\rm act}(\pi_\theta(o,\ell),a)\big],
\end{equation}
where $\alpha_k$ controls the skill mixture and $\ell_{\rm act}$ is the supervised action-fitting loss. In the residual implementation, fitting uses action regression; it need not update every VLA parameter. The policy is subsequently evaluated on each retained skill, since low fitting loss alone does not establish retention.

Historical data also include scripted/oracle-assisted experience, which must be identified separately from policy self-imitation. \pending{Specify dense-to-binary switch points, binary RL/SI ordering, bank sizes, mixing weights, and selected checkpoints for each skill.} These details must follow the actual training record rather than a retrospective assumption that every stage used an identical schedule.

\subsection{Staged RL for Complete-Task Execution}
The consolidated skills initialize learning of the complete instruction from ordinary empty-gripper starts. We combine skills by training one policy on the full task with milestone feedback. The final policy must generate the whole sequence under the final instruction; success is not obtained by externally switching between separately prompted specialists. Teacher chains may supply initialization experience, but full-task evaluation does not use teacher assistance.

Milestones credit progress through approach, pickup, and placement while retaining strict task completion as the endpoint. This exposes composition failures that separate skill tests cannot reveal. The complete-task CDPR stage uses sparse milestones, distinct from the dense reward used for initial skill discovery.

Two choices determine how intermediate feedback enters learning: which progress is rewarded, and which actions receive its credit. We implement phase-aware dense feedback for reward-construction experiments and milestone advantages for CDPR composition. The approach, held-pickup, and placement segments correspond conceptually to reach, grasp/lift, and move/place phases; the predicates and return constructions are specified for each implementation. Both distinguish meaningful partial achievements that terminal failure alone would conflate.

% Joint-paper method: phase-aware dense feedback evaluated on LIBERO.
\subsection{Phase-Aware Dense Feedback}\label{sec:dense_method}
Our dense implementation represents supported pick-and-place behavior with five ordered phases: reach, grasp, lift, move, and place. Each environment maintains its own tracker from simulator state. If $p_t\in[0,1]$ denotes active-phase progress and $B_t$ the best progress previously reached in that phase, the rewarded increment is $\Delta p_t^+=\max(p_t-B_t,0)$. The baseline resets on phase advancement. Revisiting an already reached state within a phase earns no additional progress credit.

Let $c_t$ count newly completed phases and $y_t$ denote terminal verifier success. The additive reward is
\begin{equation}
r_t=\lambda y_t+\operatorname{clip}(w_p\Delta p_t^++w_c c_t,-C,C).
\label{eq:dense}
\end{equation}
Only shaping is clipped; the original terminal reward is retained without a second terminal term. Setting $w_p=0$ gives phase-only shaping: intermediate credit occurs at phase boundaries. With $w_p>0$, new progress can earn credit before the next boundary. Thus phase awareness and reward density are separate choices. Both retain achievement information, while our CDPR implementation additionally controls stage-specific credit. \pending{Insert exact phase predicates, progress functions, transition/reset logic, and chunk-to-reward aggregation from source code.}

Rollout selection is a separate mechanism. A group of $n$ trajectories with terminal outcomes $Y_i$ is retained when $L\le n^{-1}\sum_iY_i\le U$. Bounds excluding zero reject all-failure groups even if dense returns differ. Dense feedback therefore does not automatically remove this filtering bottleneck. Both the reward law and the acceptance rule must be reported when interpreting update throughput.


\subsection{Action Module in the SmolVLA Implementation}
A frozen SmolVLA model produces a mapped prior chunk $p$. The residual input $x$ combines proprioception and frozen visual features; the residual network also receives $p$. Its bounded mean is
\begin{equation}
\mu_\theta(x,p)=\tanh\!\left(p+s_r\tanh f_\theta(x,p)\right).
\label{eq:actor}
\end{equation}
Both nonlinearities are retained: $s_r$ bounds the correction before the outer action squashing. The inspected configuration uses $s_r=1$, eight-slot chunks with replanning every four executed slots, six proprioceptive dimensions, and a 512-dimensional frozen visual projection. The long continuation trains the residual and exploration parameters while keeping inherited VLA/LoRA weights frozen. \pending{Confirm these dimensions and controller frequency from the selected run's resolved configuration.}

\subsection{Group-Relative Updates with Stage Credit}
The full-task collector records three binary milestones: approach, pickup, and strict placement, denoted $M_{i1},M_{i2},M_{i3}=Y_i$ for candidate $i$. Within each scene group, milestone returns are centered and divided by their population standard deviation, with numerical flooring and configured advantage clipping. Let $A_{ik}$ denote this advantage and $U_{ik}$ the corresponding usability mask. Constant-return groups provide no milestone contrast. Invalid candidates are excluded from training statistics, although numerical failures remain failures in evaluation.

For approach and pickup records, a terminal bonus supplements local credit:
\begin{align}
\widetilde A_{ik}&=U_{ik}A_{ik}+\beta U_{i3}T_\kappa(A_{i3},M_{ik}),\label{eq:credit}\\
T_\kappa(a,m)&=\begin{cases}\kappa a,&m=1,\ a<0,\\a,&\text{otherwise}.\end{cases}
\end{align}
Placement uses $U_{i3}A_{i3}$ without an additional duplicate bonus. The inspected settings use $\beta=1$ and $\kappa=0$: a failed completion does not contribute a negative terminal bonus to an achieved earlier milestone. This protects one credit component, not every parameter supporting the prefix.

With fixed batch advantages, nonnegative stage/candidate weights $w_r$, and per-record likelihood ratios $\rho_r$, the clipped surrogate is
\begin{equation}
S(\theta)=\sum_r w_r\min\{\rho_r\widetilde A_r,
\operatorname{clip}(\rho_r,1-\epsilon_-,1+\epsilon_+)\widetilde A_r\}.
\end{equation}
The weights balance retained stage mass and candidate trajectory lengths. Entropy and mean-action penalties accompany the surrogate, and Adam and gradient clipping determine the actual update. Thus the implemented objective is more specific than the sum of three milestone rewards.

\subsection{Conditional Latent-Action Likelihood}
At each sampled slot, exploration produces
\begin{equation}
u_t=\mu_\theta(x_t,p_t)+b_t+\sigma_\theta\odot\varepsilon_t,
\qquad a_t=\operatorname{clip}(u_t,-1,1),
\end{equation}
where $\varepsilon_t$ is standard Gaussian and $b_t$ is a gated contribution from an exogenous episode-level offset. The collector retains the unclipped sample $u_t$, the realized offset, and the old log probability. New/old ratios score the same sample under the conditional Gaussian density $\mathcal N(\mu_\theta+b_t,\operatorname{diag}\sigma_\theta^2)$.

Scoring the clipped action as a Gaussian observation incorrectly represents probability mass at the bounds. Treating a shared offset as fresh independent noise at every step also loses temporal dependence. Conditioning on its recorded realization avoids both mismatches. It does not eliminate biases introduced by stage selection, invalid-candidate filtering, or repeated surrogate optimization.

